\documentclass[11pt,a4paper]{article}

\usepackage[T1]{fontenc}
\usepackage[utf8]{inputenc}
\usepackage[spanish]{babel}
\usepackage[a4paper,margin=2.8cm]{geometry}
\usepackage{microtype}
\usepackage{xurl}
\usepackage[hidelinks]{hyperref}

\hypersetup{
  pdftitle={Reconstruir vínculos: un comentario},
  pdfauthor={Juan José Gómez Cadenas}
}

\setlength{\parindent}{1.25em}
\setlength{\parskip}{0.35em}
\raggedbottom

\title{Reconstruir vínculos: un comentario}
\author{Juan José Gómez Cadenas}
\date{}

\begin{document}

\maketitle

Recuerdo aún unos versos del gran poeta cubano José Martí que aprendí en mi juventud, escuchando al querido y añorado Pablo Milanés:

\begin{verse}
Hay una raza vil de hombres tenaces\\
De sí propios inflados, y hechos todos,\\
Todos del pelo al pie, de garra y diente;\\
Y hay otros, como flor, que al viento exhalan\\
En el amor del hombre su perfume.
\end{verse}

Vivimos en un tiempo que se diría cautivo de esa raza vil. Y el común de los mortales, frente a su sostenida agresión, acaba por resentirse, por descreer, por perder la fe. Malo es que las iglesias estén vacías, pero aún peor es que también se vacíe el ágora donde se debate lo que es de todos. Los otrora ciudadanos abandonan y se rinden, desertan a la tierra de nadie del voto en blanco o se suman a las trincheras de uno u otro bando. El resultado es una convivencia herida que se rompe un poco más cada día.

Leer un manifiesto que empieza por pedir que se reconstruyan vínculos suena extraño. ¿Qué quieren estos de mí?, se pregunta uno y la respuesta asusta un poco.

Dicen querer la amistad cívica.

Dan voz a Felipe VI, un rey que insiste en la necesidad de encontrarnos y trabajar juntos.

Y a León XIV, un papa que se atreve a llamar magnífica a la humanidad y a hablar de la cultura del amor.

Dicen:

\begin{quotation}
Queremos participar en la vida política sin ceder al escepticismo y sin depositar en ella toda nuestra esperanza.
\end{quotation}

Dicen:

\begin{quotation}
Por esta razón, nos parece urgente poner en el centro de la política la persona humana, el horizonte del bien común y una visión objetiva del papel que cumple la política en la sociedad.
\end{quotation}

Dicen:

\begin{quotation}
el voto expresa siempre una opción contingente y no supone una adhesión acrítica a un determinado partido, ni un cheque en blanco ni definitivo.
\end{quotation}

Piden:

\begin{quotation}
una fórmula de gobierno que sostenga el diálogo, superar la lógica de la contraposición, derribar muros, respetar la libertad, renunciar al adoctrinamiento, regenerar las instituciones y trabajar juntos.
\end{quotation}

Solo unos insensatos se atreven a tender la mano cuando hay tantos dispuestos a cortársela.

El mundo y este país necesitan, desesperadamente, más insensatos como estos.

Me recuerdan los versos de mi juventud, a esos otros hombres en los que hay que creer:

\begin{verse}
como flor, que al viento exhalan\\
En el amor del hombre su perfume.
\end{verse}

\end{document}
